\chapter{相关研究与理论基础}

\section{职业生命周期视角}

职业发展理论强调个体会在探索、建立、维持和再适应之间不断移动。当媒介环境发生变化时，这一过程会被进一步放大，因为既有的声誉资本并不会自动转化为新平台中的竞争优势。

\section{形象管理视角}

自我呈现理论指出，公众人物需要在不同场景中组织可接受的角色形象。社交平台进入职业系统后，这种呈现从阶段性展示转向持续性协商，危机也更容易被长期标签化。由此可见，形象韧性并不只是“重新出现”，而是“重新被理解”。

\section{平台选择视角}

不同平台在可达性、互动性和自主性方面存在显著差异。传统媒体有助于建立集中曝光，影视与综艺有助于扩大角色边界，而视频平台则更有利于重新掌控表达节奏。平台选择因此成为观察职业转型的重要切口。
